% -----------------------------------------------------------------------
\section{Results}
\label{sec:results}
% -----------------------------------------------------------------------

We evaluate the surrogate on two cases: the idealised Xie--Castro building-array
benchmark and the realistic Barcelona urban block. For each case we report
inference (autonomous rollout), uncertainty quantification, and data
assimilation.

% =======================================================================
\subsection{Xie--Castro}
\label{sec:xie_castro}
% =======================================================================

\textit{This case is under preparation.} The Xie--Castro configuration is an
idealised, regular building-array benchmark used to validate the surrogate under
controlled conditions before moving to the full Barcelona geometry.

\subsubsection{Inference}
\label{sec:xc_inference}

% TODO: Xie--Castro autonomous rollout evaluation.
\textit{Placeholder for the Xie--Castro rollout evaluation.}

\subsubsection{Uncertainty Quantification}
\label{sec:xc_uq}

% TODO: Xie--Castro uncertainty quantification.
\textit{Placeholder for the Xie--Castro uncertainty quantification.}

\subsubsection{Data Assimilation}
\label{sec:xc_da}

% TODO: Xie--Castro ES-MDA experiment.
\textit{Placeholder for the Xie--Castro data-assimilation experiment.}

% =======================================================================
\subsection{Barcelona}
\label{sec:barcelona}
% =======================================================================

\subsubsection{Inference}
\label{sec:bcn_inference}

\begin{figure}[t]
  \centering
  \includegraphics[width=\textwidth]{rollout.png}
  \caption{Autonomous surrogate rollout on a held-out Barcelona trajectory.
    Each column shows a snapshot of the velocity magnitude
    $\umag = \sqrt{u^2+v^2+w^2}$ at ground level ($z$-index 0) at five
    equally spaced instants covering the full 120-step rollout.
    \textbf{Top:} \palm{} ground truth.
    \textbf{Middle:} \pthreeds{} surrogate prediction, initialised from the
    truth state at $t=0$ and then rolled out autoregressively.
    \textbf{Bottom:} absolute error $|\,|\hat{\mathbf{u}}| - |\mathbf{u}|\,|$
    (note the separate, saturated colour scale).
    The surrogate reproduces the large-scale flow organisation and street-canyon
    patterns throughout the rollout; error concentrates at building edges and
    in the more energetic street-level jets.}
  \label{fig:rollout}
\end{figure}

\begin{figure}[t]
  \centering
  \begin{subfigure}[b]{0.48\textwidth}
    \includegraphics[width=\textwidth]{rmse.png}
    \caption{Per-step rollout RMSE (normalised velocity units).
      Error accumulates steeply over the first $\sim$20 steps as the
      surrogate diverges from the exact phase trajectory, then saturates
      near the climatological error level as the rollout reaches a
      statistically stationary regime.}
    \label{fig:rmse}
  \end{subfigure}
  \hfill
  \begin{subfigure}[b]{0.48\textwidth}
    \includegraphics[width=\textwidth]{params.png}
    \caption{Inflow parameters for the evaluation rollout:
      inflow angle $\iang$ (top) and velocity magnitude $\ivel$ (bottom)
      as functions of rollout time step.
      Both follow smooth AR(2) trajectories spanning a moderate range
      of wind directions and speeds representative of the Barcelona site.}
    \label{fig:params}
  \end{subfigure}
  \caption{Quantitative rollout evaluation. (\subref{fig:rmse}) RMSE profile
    over 120 autonomous steps. (\subref{fig:params}) Inflow parameter
    trajectory used for the evaluation rollout.}
  \label{fig:quant}
\end{figure}

Figure~\ref{fig:rollout} illustrates the surrogate's ability to propagate
the Barcelona urban-block flow field over 120 autonomous time steps under
smoothly varying inflow conditions (Figure~\ref{fig:params}).
The predicted fields preserve the main flow structures — accelerated jets
in street canyons, recirculation zones in building wakes, and the
roof-level shear layer — throughout the rollout. Errors are largest at
building edges and in regions of strong velocity gradients, consistent
with the expected difficulty of predicting the exact phase of turbulent
detachment.

Figure~\ref{fig:rmse} shows that the per-step RMSE (in normalised velocity
units) rises steeply during the first $\sim$20 steps as the surrogate
rolls out from a precise initial condition into its own trajectory, before
saturating near a climatological error floor around step 30--40. This
saturation indicates that the surrogate orbit remains physically plausible
and bounded — a necessary property for stable ensemble data assimilation,
where individual members are not expected to phase-lock to the truth but
must collectively span the plausible state space.

\subsubsection{Uncertainty Quantification}
\label{sec:bcn_uq}

% TODO: Barcelona uncertainty quantification.
\textit{Placeholder for the Barcelona uncertainty quantification.}

\subsubsection{Data Assimilation}
\label{sec:da_results}

We now embed the surrogate in the \esmda{} loop of Section~\ref{sec:esmda} and
assimilate synthetic street-level observations drawn from a held-out \palm{} truth
trajectory. The truth is a $1200\,\mathrm{s}$ \palm{}-LES integration of the
Barcelona block sampled at $\Delta t = 10\,\mathrm{s}$ ($120$ frames), generated
by the full LES solver and \emph{not} seen during surrogate training, so the
experiment is free of the inverse crime (Section~\ref{sec:warmstart}). The
forward model in the loop is the trained \pthreeds{} surrogate. An ensemble of
$N = 50$ members is propagated over $4$ sequential assimilation windows of
$300\,\mathrm{s}$ ($30$ surrogate steps) each, with $\niter = 3$
\esmda{} analysis passes per window, observation-error standard deviation
$\sigma_o = 0.25$ (velocity units), and no covariance localisation or inflation.
Observations are the temporally averaged velocity components at the six
street-level sensors of Section~\ref{sec:obs}. The estimated control variable is
the time-varying inflow-parameter trajectory $(\iang, \ivel)$ (the
\texttt{TimeVaryingParameterESMDA} smoother); the velocity field is propagated by
the surrogate conditioned on each member's parameters and warm-start initial
state, and each window's posterior seeds the next.

\paragraph{Inflow-parameter estimation.}
Figure~\ref{fig:da_param_evolution} shows the prior and posterior inflow-parameter
ensembles across the four windows. The broad AR(2) prior (orange) collapses onto a
tight posterior (blue) that tracks the truth (dashed) closely for the inflow
\emph{angle}: the posterior follows the slow swings of the true wind direction,
and the per-step angle RMSE (Figure~\ref{fig:da_param_error}, top) falls from a
prior mean of $11.1^\circ$ to a posterior mean of $7.3^\circ$ (final-time
$6.3^\circ$), a $34\%$ reduction, with a continuous-ranked-probability score
(CRPS) of $6.0^\circ$. The inflow \emph{velocity magnitude} is harder: the
posterior is pulled systematically high (settling near $8.5$--$9.5\,\mathrm{m/s}$
against a true $\sim\!7.5\,\mathrm{m/s}$) and is over-confident, so its
time-averaged RMSE ($1.41\,\mathrm{m/s}$) actually \emph{exceeds} the prior
($0.95\,\mathrm{m/s}$). The estimate does improve window-on-window, recovering to
a final-time RMSE of $0.88\,\mathrm{m/s}$ (just below the prior level), but the
early-window overshoot dominates the average. This degradation is an
observability limitation rather than a failure of the update: the surrogate's
ground-level speeds at the sensor sites are smoother and lower-amplitude than
\palm{}'s, so fitting the sensor data biases the inferred inflow speed upward.

\begin{figure}[t]
  \centering
  \includegraphics[width=\textwidth]{esmda_param_evolution.png}
  \caption{Inflow-parameter evolution over the four assimilation windows
    (shaded/unshaded bands). \textbf{Top/middle:} prior ensemble (orange) and
    posterior ensemble (blue) for the inflow angle and velocity magnitude, with
    the \palm{} truth dashed. The posterior tightens sharply relative to the prior;
    the angle tracks the truth, while the speed is biased high.
    \textbf{Bottom:} domain-averaged velocity-magnitude RMSE of the posterior
    state over the full $120$-step horizon, which stays bounded across all four
    windows (no rollout divergence under assimilation).}
  \label{fig:da_param_evolution}
\end{figure}

\begin{figure}[t]
  \centering
  \includegraphics[width=0.85\textwidth]{esmda_param_error.png}
  \caption{Per-step posterior error for the two inflow parameters: RMSE (blue)
    and CRPS (orange) of the inflow angle (top) and velocity magnitude (bottom)
    over the four windows. The angle error is well below its prior level
    throughout; the velocity-magnitude error is largest in the first window and
    relaxes thereafter.}
  \label{fig:da_param_error}
\end{figure}

\paragraph{State reconstruction.}
Figure~\ref{fig:da_final_state} compares the truth, posterior-mean, and
posterior-standard-deviation velocity-magnitude fields at the final time. The
posterior mean reproduces the large-scale organisation of the truth---the
street-canyon channelling, the building-wake voids, and the higher-energy region
along the eastern (downstream) edge of the block. The ensemble standard deviation
is smallest within the densely built core, where the geometry constrains the flow,
and concentrates along the main channels and in the energetic downstream region
that the six sensors do not cover---an intuitive map of where the remaining
posterior uncertainty lives. The domain-averaged state RMSE is $1.40\,\mathrm{m/s}$
(min $0.74$, max $1.91\,\mathrm{m/s}$) and, as the bottom panel of
Figure~\ref{fig:da_param_evolution} shows, remains bounded over the entire
$1200\,\mathrm{s}$---confirming that the surrogate rollout stays stable inside the
repeatedly re-initialised assimilation loop.

\begin{figure}[t]
  \centering
  \includegraphics[width=\textwidth]{esmda_final_state.png}
  \caption{Velocity-magnitude field at the final assimilation time.
    \textbf{Left:} \palm{} truth, with the six assimilation sensors marked.
    \textbf{Centre:} \esmda{} posterior-mean. \textbf{Right:} posterior ensemble
    standard deviation. The posterior mean captures the truth's canyon and wake
    structure; spread concentrates along the main channels and the unobserved
    downstream region.}
  \label{fig:da_final_state}
\end{figure}

\paragraph{Sensor fit and observability.}
Figure~\ref{fig:da_sensors} overlays the ensemble, its mean, and the truth at the
six sensors. The posterior ensemble brackets the truth on average but cannot
reproduce its sharp gusts: the \palm{} time series at street level are considerably
more intermittent than the surrogate's smoother ground-level dynamics, and with
only six point sensors the high-frequency content is effectively unobservable.
The mean sensor velocity-vector RMSE is $0.80\,\mathrm{m/s}$ with an energy score
of $0.59$. This under-representation of street-level variance is the same effect
that biases the inferred inflow speed, and points to denser or better-placed
sensing---rather than a larger ensemble---as the route to further improvement.

\begin{figure}[p]
  \centering
  \includegraphics[width=0.82\textwidth]{esmda_sensors.png}
  \caption{Velocity magnitude $\Umag$ at the six street-level
    assimilation sensors ($z = 3\,\mathrm{m}$): posterior ensemble members (thin
    blue), ensemble mean (thick blue), and \palm{} truth (dashed). The bottom panel
    is the across-sensor RMSE and CRPS. The ensemble tracks the slow trends but
    misses the truth's sharp gusts, reflecting the limited observability of
    street-level turbulence from six points.}
  \label{fig:da_sensors}
\end{figure}

\paragraph{Cost.}
The complete assimilation---$50$ members $\times$ $4$ windows $\times$ $3$ \esmda{}
passes, i.e.\ $200$ surrogate rollouts of $30$ steps---completed in
$5.4\times10^{3}\,\mathrm{s}$ ($\approx 1.5\,\mathrm{h}$) of wall time on a single
GPU, averaging $23\,\mathrm{min}$ per window. The same ensemble run with the \palm{}
solver in the loop would require $200$ LES integrations of the Barcelona domain,
orders of magnitude more expensive, which is precisely the cost the surrogate is
designed to amortise.

Table~\ref{tab:da_metrics} collects the headline metrics. The overall picture is
that the surrogate-driven \esmda{} recovers the dominant inflow direction and a
spatially faithful posterior mean at a fraction of the LES cost, while the
street-level velocity amplitude remains limited by sensor observability---a
configuration issue (sensor density and placement) rather than a deficiency of
the surrogate or the update scheme.

\begin{table}[t]
  \centering
  \caption{Headline \esmda{} metrics ($50$ members, $4$ windows, $3$ passes;
    \palm{} truth, \pthreeds{} forward model). Parameter and state errors are averaged
    over the full $1200\,\mathrm{s}$ horizon; ``prior'' is the corresponding
    pre-assimilation ensemble error.}
  \label{tab:da_metrics}
  \renewcommand{\arraystretch}{1.15}
  \footnotesize
  \begin{tabular}{llll}
    \toprule
    Quantity & Prior RMSE & Posterior RMSE & CRPS \\
    \midrule
    Inflow angle $\iang$ (deg)        & $11.1$ & $7.3$ \,(final $6.3$; $-34\%$) & $6.0$ \\
    Inflow velocity $\ivel$ (m\,s$^{-1}$)  & $0.95$ & $1.41$\,(final $0.88$; $+49\%$) & $1.27$ \\
    State $\Umag$ (m\,s$^{-1}$) & ---    & $1.40$\,(min $0.74$, max $1.91$) & --- \\
    Sensor velocity, vector (m\,s$^{-1}$)& ---    & $0.80$ & $0.59$ (ES) \\
    \midrule
    \multicolumn{4}{l}{Wall time: $5.4\times10^{3}\,\mathrm{s}$ total
      ($\approx 23\,\mathrm{min}$/window), single GPU.} \\
    \bottomrule
  \end{tabular}
\end{table}
